\section{Introduction}

Ce laboratoire porte sur l'étude d'un circuit magnétique en double E, avec un entrefer réglable, une bobine primaire (\(N_1=500\) spires) et une bobine secondaire de mesure (\(N_2=200\) spires). L'objectif est de relier les modèles théoriques (résistance, flux, induction, inductances, couplage AC) aux observations expérimentales.

Les buts sont :
\begin{itemize}
\item établir les expressions analytiques des grandeurs magnétiques dans deux hypothèses (\(\mu_r\to\infty\) et \(\mu_r=400\));
\item tracer les courbes théoriques demandées par l'énoncé ;
\item comparer les mesures (résistance, \(B=f(I)\), diagramme \(x-y\) et intégrateur) à la théorie ;
\item discuter la validité des hypothèses (franges, saturation, non-linéarités, pertes).
\end{itemize}

Moyens utilisés : maquette magnétique, multimètre (2 fils / 4 fils), teslamètre ou sonde Hall, générateur de fonctions, amplificateur linéaire, sonde de courant, oscilloscope et intégrateur à ampli-op.
